\documentclass[10pt,a4paper]{amsart}
\usepackage[margin=19mm]{geometry}
\usepackage{amsmath,amssymb,mathtools}
\usepackage[scr=rsfs,cal=euler]{mathalfa}
\usepackage{xfrac}
\usepackage[unicode,hidelinks,bookmarks=false]{hyperref}
\newcommand{\ie}{i.e.\ }
\newcommand{\cf}{cf.\ }
\newcommand{\resp}{respectively\ }
\newcommand{\Subsection}{\subsection}
\providecommand{\nobreakdash}{\nobreak\hbox{-}}
\setlength{\emergencystretch}{2em}
\multlinegap0pt
\usepackage{tikz}
\usepackage{amssymb}
\usepackage{float}
\newtheorem{lem}{Lemma}
\newcommand{\C}{\mathbb{C}}
\newcommand{\R}{\mathbb{R}}
\title{Characterisation of some multivalued harmonic functions on $\R^{n}$}
\author{Yang Li, Parsa Mashayekhi, Yichen Zhang}
\date{Source: arXiv:2608.15722v1 (CC BY 4.0)}
\begin{document}
\maketitle

\noindent\emph{Source and licence.} Characterisation of some multivalued harmonic functions on $\R^{n}$. Version arXiv:2608.15722v1 (CC BY 4.0), distributed by arXiv under the Creative Commons Attribution 4.0 International licence, http://creativecommons.org/licenses/by/4.0/, as recorded in the arXiv licence metadata for this exact version and shown on the abstract page https://arxiv.org/abs/2608.15722v1. Full licence notice: \texttt{licenses/arxiv-2608.15722-CC-BY-4.0.txt}.\par\smallskip

\noindent\textit{Editorial scope.} Counted unit: the article's lem lem:floerdegreeforgraph (source0.tex lines 855--859). The article's complete original proof of it is delivered with it (source0.tex lines 860--910, one \texttt{proof} environment of 51 lines). No mathematics is written, repaired or re-derived: the statement and the proof are the article's own lines, in the article's own notation, and no retained line is substituted or re-flowed.  No further source-local context is needed: every cross-reference of every retained line resolves inside the two delivered spans.  Nothing was omitted from any retained span, so the package carries no omission record.  No retained line contains a citation, so the wrapper carries no bibliography and no bibliography entry.  No retained line contains a figure environment, an image-inclusion command or drawing code, so no image is delivered and the package declares no asset dependency.  Editorial formatting only, in addition: an AMS \texttt{amsart} preamble that restates the article's own declarations and macros used by the retained lines, namely the environments \texttt{lem} on separate counters, together with the standard packages those lines need.  The retained environments print in this document as Lemma 1 in order of appearance, whereas the article prints the same items with the numbering of its own full document.  The complete native source of the article as distributed in the arXiv e-print for this exact version is bundled unchanged as \texttt{sources/207776/source0.tex}.
\begin{lem}\label{lem:floerdegreeforgraph} 
      Let $F_1,F_2$ be two (single-valued) $C^\infty$-functions defined near $0\in \R^n$, with $|\nabla^2 F_1|,|\nabla^2 F_2|\leqslant 1$. Suppose that $dF_1(0)=dF_2(0)=p$, and $\nabla^2(F_2-F_1)(0)$ is non-degenerate. Then
      for $0<\tau\leqslant \tau_0(n)\ll 1$, the Floer degree at $(0,\tau p)$ agrees with the negative index of $\nabla^2(F_2-F_1)(0)$, 
      $$\mu_{L_{\tau dF_1},L_{\tau dF_2}}(0,\tau p)=\operatorname{ind}(\nabla^2(F_2-F_1)(0)),$$ where the Lagrangian angles of $L_{\tau dF_i}$ are defined to be $\theta_i:=\sum\limits_{k=1}^n \arctan (\tau \lambda_k^i)$ for the eigenvalues $\lambda_1^i,\cdots,\lambda_n^i$ of $\nabla^2F_i$ near $0$, and $i=1,2$.    
   \end{lem}

   \begin{proof} The main effect of compressing two graphs by $\tau\leqslant \tau_0\ll1$ is that we can freely use the identity 
   $$\arg \det(I+\sqrt{-1}\tau A)=\sum_{k=1}^n\arctan (\tau \mu_k), $$
for the eigenvalues $\mu_1,\cdots,\mu_n$ of a real symmetric matrix $\|A\|\lesssim 1$, and $\arg :\C\setminus \R_{<0}\to (-\pi,\pi)$.  
   
       Writing $A_i=\nabla^2F_i(0)$, the tangent plane at $(0,\tau p)$ is given by
       $$T_i:=T_{(0,\tau p)}L_{\tau dF_i}=\operatorname{Graph}(\tau\nabla^2F_i(0))=(I+\sqrt{-1}\tau A_i)\R^n,\quad i=1,2, $$
       where the origin of $\R^n$ corresponds to $(0,\tau p)$ in the original coordinates. 
       
       Then $U_1:=(I+\tau^2 A_1^2)^{-1/2}(I-i\tau A_1)\in U(n)$ satisfies
       $$U_1T_1=\R^n,\quad |U_1-Id|\lesssim \tau|A_1|\leqslant\tau,\quad \arg \det U_1=-\theta_{1}(0).$$
       To write $U_1T_2$ as a graph over $\R^n$, we calculate that 
       $$U_1T_2=(I+\tau^2A_1^2)^{-1/2}(I+\tau^2A_1A_2+\sqrt{-1}\tau (A_2-A_1))\R^n=(X_{\tau}+\sqrt{-1}\tau Y_\tau)\R^n,$$
       where $X_\tau=(I+\tau^2A_1^2)^{-1/2}(I+\tau^2A_1A_2)$ is invertible provided that $\tau\leqslant \frac{1}{2}$ and $Y_\tau=(I+\tau^2A_1^2)^{-1/2}(A_2-A_1)$. Hence $U_1T_2=\operatorname{Graph}(\tau Y_{\tau}X_{\tau}^{-1})$ over $U_1T_1=\R^n$. Since $A_2-A_1$ is non-degenerate, so is $Y_{\tau}X_{\tau}^{-1}$. We compute 
        \begin{equation}
        \begin{aligned}
            \label{eqn:symmetricgraphmatrix}
         X_\tau^TY_\tau&=(I+\tau^2A_2A_1)(I+\tau^{2}A_1^{2})^{-1}(A_2-A_1)
               \\&=(I+\tau^2(A_1+A_2-A_1)A_1)(I+\tau^2A_1^2)^{-1}(A_2-A_1)
               \\&=(A_2-A_1)+\tau^2(A_2-A_1)A_1(I+\tau^2A_1^2)^{-1}(A_2-A_1),
            \end{aligned}
        \end{equation}
       and 
       \begin{equation}
           \begin{aligned}
               Y_\tau^T X_\tau&=(A_2-A_1)(I+\tau^2A_1^2)^{-1}(I+\tau^2A_1A_2)
               \\&=(A_2-A_1)(I+\tau^2A_1^2)^{-1}(I+\tau^2A_1(A_1+A_2-A_1))
               \\&=(A_2-A_1)+\tau^2(A_2-A_1)(I+\tau^2A_1^2)^{-1}A_1(A_2-A_1)
               \\&=(A_2-A_1)+\tau^2(A_2-A_1)A_1(I+\tau^2A_1^2)^{-1}(A_2-A_1)\quad (A_1,(I+\tau^2A_1^2)^{-1} \textnormal{commute})
               \\&=X_\tau^TY_\tau.
           \end{aligned}
       \end{equation}
       It follows that $Y_\tau X_\tau^{-1}$ is symmetric. 
       Now, writing $\tau \mu_1,\cdots,\tau\mu_n$ as the eigenvalues of $\tau Y_\tau X_\tau^{-1}$ with respect to an orthogonal basis of $U_1T_1=\R^n$, we have 
       $$U_1T_2=(e^{i\alpha_1},\cdots,e^{i\alpha_n})\R^n,\quad \alpha_k=\left\{\begin{array}{ll}
          \arctan(\tau\mu_k),  & \mu_k>0, \\
            \arctan(\tau\mu_k)+\pi, & \mu_k<0,
       \end{array}\right.\quad 1\leqslant k\leqslant n.$$
       Thus, 
       $$\mu_{L_{\tau dF_1},L_{\tau dF_2}}(0,\tau p)=\frac{1}{\pi}\left(\sum_{k=1}^n \alpha_k+\theta_1(0)-\theta_2(0)\right)=\operatorname{ind}(Y_\tau X_{\tau}^{-1}),$$
       as $\theta_2(0)-\theta_1(0)=\arg\det (U_1T_2)=\arg\left(\det(I+\sqrt{-1}\tau Y_\tau X_\tau^{-1})\det(X_\tau)\right)=\sum\limits_{k=1}^n\arctan(\tau\mu_k)$, whenever $\tau\leqslant \tau_0\ll1$. It remains to show 
       $$\operatorname{ind}(Y_{\tau}X_{\tau}^{-1})=\operatorname{ind}(A_2-A_1).$$
       Indeed, for $0<\tau\leqslant\tau_0(n)\ll1$, as the index is preserved by the congruence operation, we obtain \begin{equation*}
           \begin{aligned}
               \operatorname{ind}(Y_{\tau}X_{\tau}^{-1})&=\operatorname{ind}(X_\tau^T Y_\tau)\quad \quad\quad (\textnormal{congruence by} ~ X_\tau)
               \\&=\operatorname{ind}((A_2-A_1)+\tau^2(A_2-A_1)A_1(I+\tau^2A_1^2)^{-1}(A_2-A_1))\quad \textnormal{by (\ref{eqn:symmetricgraphmatrix})}
               \\&=\operatorname{ind}(\underbrace{(A_2-A_1)^{-1}}_{\textnormal{norm of all eigenvalues}\geqslant C(n)} +\underbrace{\tau^2A_1(I+\tau^2A_1^2)^{-1})}_{\textnormal{matrix norm}\leqslant C\tau^2}\quad (\textnormal{congruence by} ~ (A_2-A_1)^{-1})
                \\&= \operatorname{ind}((A_2-A_1)^{-1})=\operatorname{ind}(A_2-A_1).
           \end{aligned}
       \end{equation*}
The Lemma follows.        
   \end{proof}

\end{document}
